\documentclass[12pt,a4paper]{report}

% ---------------------------------------------------------------------------
% SkillBridge academic report
% ---------------------------------------------------------------------------
\usepackage[utf8]{inputenc}
\usepackage[T1]{fontenc}
\usepackage[english]{babel}
\usepackage{geometry}
\usepackage{graphicx}
\usepackage[table]{xcolor}
\usepackage{sectsty}
\usepackage{hyperref}
\usepackage{booktabs}
\usepackage{longtable}
\usepackage{tabularx}
\usepackage{array}
\usepackage{enumitem}
\usepackage{float}
\usepackage{caption}
\usepackage{subcaption}
\usepackage{listings}
\usepackage{pdflscape}

\geometry{top=2cm,bottom=2cm,left=2cm,right=2cm}
\setlength{\parindent}{0pt}
\setlength{\parskip}{0.45em}
\renewcommand{\arraystretch}{1.2}
\hypersetup{colorlinks=true,linkcolor=sbNavy,urlcolor=sbTeal,citecolor=sbNavy}

\definecolor{sbNavy}{HTML}{1F3A5F}
\definecolor{sbTeal}{HTML}{1B998B}
\definecolor{sbOrange}{HTML}{E67E22}
\definecolor{sbSky}{HTML}{4C86A8}
\definecolor{sbGold}{HTML}{C9A227}
\definecolor{sbMint}{HTML}{6CBFB4}
\definecolor{sbGray}{HTML}{F4F6F8}
\definecolor{sbDark}{HTML}{20242A}
\definecolor{sbLine}{HTML}{D8DEE9}
\arrayrulecolor{sbLine}

\chapterfont{\color{sbNavy}}
\sectionfont{\color{sbDark}}
\subsectionfont{\color{sbTeal}}

\lstdefinestyle{reportcode}{
    basicstyle=\ttfamily\small,
    breaklines=true,
    frame=single,
    rulecolor=\color{sbLine},
    backgroundcolor=\color{sbGray},
    columns=fullflexible,
    keepspaces=true,
    showstringspaces=false
}
\lstset{style=reportcode}

\newcommand{\projectName}{SkillBridge}
\newif\ifimagefound
\let\originalfbox\fbox
\renewcommand{\fbox}[1]{%
    \ifimagefound
        \global\imagefoundfalse
    \else
        \originalfbox{#1}%
    \fi
}
\newcommand{\maybeincludegraphics}[2][]{
    \IfFileExists{#2}
        {\global\imagefoundtrue\includegraphics[#1]{#2}\par\vspace{0.35cm}}
        {\global\imagefoundfalse}
}
\newcommand{\diagramimage}[3]{
    \begin{figure}[H]
        \centering
        \IfFileExists{#1}
        {\includegraphics[width=#2\textwidth]{#1}}
        {
            \fcolorbox{sbLine}{sbGray}{
                \begin{minipage}[c][7cm][c]{0.88\textwidth}
                    \centering
                    {\Large Diagram image missing}\\[0.5cm]
                    \texttt{\detokenize{#1}}\\[0.35cm]
                    Render the matching PlantUML file from \texttt{plantuml/} and save the image here.
                \end{minipage}
            }
        }
        \caption{#3}
    \end{figure}
}
\newcommand{\screenshotplaceholder}[3]{
    \begin{figure}[H]
        \centering
        \IfFileExists{#1}
        {\includegraphics[width=0.92\textwidth]{#1}}
        {
            \fcolorbox{sbLine}{sbGray}{
                \begin{minipage}[c][7cm][c]{0.88\textwidth}
                    \centering
                    {\Large #2}\\[0.5cm]
                    \texttt{\detokenize{#1}}\\[0.35cm]
                    Add the application screenshot here before compiling the final PDF.
                \end{minipage}
            }
        }
        \caption{#3}
    \end{figure}
}

\begin{document}

% ---------------------------------------------------------------------------
% Cover
% ---------------------------------------------------------------------------
% ---------------- Page de garde ----------------
\begin{titlepage}
    \thispagestyle{empty}

    \noindent
    \begin{minipage}[t]{0.48\textwidth}
        \IfFileExists{assets/bdiologo.png}
            {\includegraphics[height=2cm,width=1.6cm]{assets/bdiologo.png}}
            {\fbox{\rule{0pt}{2cm}\hspace{1cm}BDIO\hspace{1cm}}}
    \end{minipage}
    \hfill
    \begin{minipage}[t]{0.48\textwidth}
        \raggedleft
        \IfFileExists{assets/ensamlogo.png}
            {\includegraphics[height=1.6cm]{assets/ensamlogo.png}}
            {\fbox{\rule{0pt}{1.6cm}\hspace{1cm}ENSAM\hspace{1cm}}}
    \end{minipage}

    \vspace{3.2cm}

    \begin{center}
        {\small Projet JEE \& BigData}

        \vspace{0.35cm}

        {\LARGE \bfseries SkillBridge}

        \vspace{0.45cm}

        {\normalsize
        Plateforme éducative orientée projet pour la recommandation de cours
}
    \end{center}

    \vspace{4.8cm}

    \noindent
    \begin{minipage}[t]{0.42\textwidth}
        {\small Réalisé par :}

        \vspace{0.35cm}

        {\small Elkhali Omar }

        \vspace{0.25cm}

        {\small Moutik Ayoub }

        \vspace{0.25cm}

        {\small Ouzzane Jamila }
    \end{minipage}
    \hfill
    \begin{minipage}[t]{0.32\textwidth}
        {\small Encadré par :}

        \vspace{0.35cm}

        {\small Pr. EL Faquih Loubna}

        \vspace{0.25cm}

        {\small Pr. HIRCHOUA Bader}

        \vspace{0.25cm}

        {\small Pr. SEKKATE Sara}
    \end{minipage}

    \vfill

    \begin{center}
        {\small AU : 25-26}
    \end{center}

\end{titlepage}

\pagenumbering{roman}
\chapter*{Résumé}
\addcontentsline{toc}{chapter}{Résumé}
\projectName{} est une plateforme de parcours d’apprentissage conçue à partir d’un constat pratique : de nombreux apprenants savent ce qu’ils souhaitent réaliser, mais ne savent pas quelles compétences ni quels cours peuvent les aider à y parvenir. L’application permet à un utilisateur de créer un compte, de soumettre une idée de projet, de recevoir les compétences détectées, d’obtenir des recommandations de cours classées, d’enregistrer des cours et de suivre sa progression. Les administrateurs peuvent gérer les utilisateurs, le catalogue, les fournisseurs, les catégories, les compétences ainsi que les écrans de suivi de l’état du module Big Data.

Le projet est organisé autour de trois applications coordonnées : un backend Spring Boot, un frontend React et un laboratoire Big Data orienté terminal. Le backend applique les principes JEE à travers une architecture en couches, des contrôleurs REST, des entités Jakarta Persistence, des repositories, des DTO, la validation et les transactions. La couche de sécurité utilise Spring Security avec une authentification JWT sans état, le hachage des mots de passe avec BCrypt, l’intégration de la connexion OAuth, l’autorisation au niveau des méthodes, le renforcement de la configuration CORS, les en-têtes de sécurité HTTP et la protection contre les attaques par force brute. Le module Big Data illustre l’ingestion batch et streaming avec Sqoop et Flume, le stockage distribué avec HDFS, l’analyse avec Hive, le traitement avec un job Java MapReduce, l’enrichissement des données avec Python et une couche de service basée sur HBase.

Ce rapport présente l’objectif, l’architecture, le modèle de données, les algorithmes, la conception de l’API, les contrôles de sécurité, le pipeline Big Data, les diagrammes UML, la stratégie de test et le plan de démonstration de l’application complète.

\tableofcontents
\listoffigures
\listoftables
\newpage
\pagenumbering{arabic}

% ---------------------------------------------------------------------------
\chapter{Étude fonctionnelle et analyse des besoins}

\section{Présentation de SkillBridge}

SkillBridge est une plateforme éducative intelligente permettant de transformer une idée de projet en un parcours d’apprentissage structuré. Contrairement aux plateformes classiques où l’utilisateur recherche directement un cours, SkillBridge adopte une approche orientée projet.

L’utilisateur décrit le projet qu’il souhaite réaliser en langage naturel. Le système analyse ensuite cette description afin de :
\begin{itemize}
    \item détecter les compétences nécessaires ;
    \item identifier les catégories pertinentes ;
    \item sélectionner les cours adaptés ;
    \item classer les recommandations selon un score explicable.
\end{itemize}

La plateforme combine :
\begin{itemize}
    \item une application web full-stack ;
    \item un moteur de recommandation ;
    \item un pipeline Big Data basé sur l’écosystème Hadoop ;
    \item un tableau de bord analytique destiné aux administrateurs.
\end{itemize}

Le projet vise ainsi à améliorer l’orientation pédagogique des utilisateurs tout en exploitant les données générées par la plateforme à des fins analytiques.

\section{Acteurs du système}

Le système met en interaction plusieurs acteurs ayant des rôles différents.

\begin{table}[H]
    \centering
    \caption{Acteurs principaux du système}
    
    \begin{tabular}{|p{4cm}|p{9cm}|}
        \hline
        \textbf{Acteur} & \textbf{Rôle} \\
        \hline
        
        Utilisateur &
        Crée un compte, consulte le catalogue, crée des idées de projet, génère des recommandations, sauvegarde des cours et suit sa progression. \\
        
        \hline
        
        Administrateur &
        Gère les utilisateurs, les cours, les catégories, les compétences et supervise les indicateurs analytiques ainsi que le pipeline Big Data. \\
        
        \hline
        
        Pipeline Big Data &
        Collecte, traite, stocke et analyse les données générées par l’application. \\
        
        \hline
    \end{tabular}
\end{table}

\section{Besoins fonctionnels}

Les besoins fonctionnels représentent les fonctionnalités que doit assurer le système afin de répondre aux attentes des utilisateurs.

\subsection{Authentification et gestion des utilisateurs}

Le système doit permettre :
\begin{itemize}
    \item la création de comptes utilisateurs ;
    \item l’authentification sécurisée ;
    \item la connexion via email et mot de passe ;
    \item la connexion via Google et GitHub ;
    \item la gestion des rôles utilisateur et administrateur.
\end{itemize}

\subsection{Gestion du catalogue de cours}

Le système doit permettre :
\begin{itemize}
    \item la consultation du catalogue de cours ;
    \item la recherche et le filtrage des cours ;
    \item la pagination des résultats ;
    \item la gestion des catégories et fournisseurs ;
    \item l’administration des ressources pédagogiques.
\end{itemize}

\subsection{Gestion des projets}

Le système doit permettre :
\begin{itemize}
    \item la création d’idées de projet ;
    \item l’analyse des descriptions fournies ;
    \item la détection des compétences associées ;
    \item la génération de recommandations personnalisées.
\end{itemize}

\subsection{Recommandation de cours}

Le moteur de recommandation doit :
\begin{itemize}
    \item analyser les mots-clés du projet ;
    \item détecter les compétences pertinentes ;
    \item associer les catégories correspondantes ;
    \item calculer un score de pertinence ;
    \item générer une liste ordonnée de cours.
\end{itemize}

\subsection{Suivi pédagogique}

Le système doit permettre :
\begin{itemize}
    \item la sauvegarde des cours ;
    \item le suivi de progression ;
    \item l’affichage des parcours enregistrés ;
    \item le suivi des activités utilisateur.
\end{itemize}

\subsection{Fonctionnalités analytiques}

Le système doit permettre :
\begin{itemize}
    \item la collecte des événements utilisateurs ;
    \item l’analyse des données générées ;
    \item la visualisation des statistiques ;
    \item la supervision du pipeline Big Data.
\end{itemize}

\section{Besoins non fonctionnels}

En plus des fonctionnalités principales, le système doit respecter plusieurs contraintes non fonctionnelles.

\subsection{Sécurité}

La plateforme doit garantir :
\begin{itemize}
    \item la protection des données utilisateurs ;
    \item la sécurisation des accès ;
    \item le contrôle des autorisations ;
    \item la protection contre les accès non autorisés.
\end{itemize}

\subsection{Performance}

Le système doit :
\begin{itemize}
    \item répondre rapidement aux requêtes utilisateurs ;
    \item séparer les traitements temps réel et analytiques ;
    \item limiter l’impact du pipeline Big Data sur l’application web.
\end{itemize}

\subsection{Scalabilité}

L’architecture doit permettre :
\begin{itemize}
    \item l’augmentation du volume de données ;
    \item l’ajout de nouveaux nœuds Hadoop ;
    \item l’extension des fonctionnalités analytiques.
\end{itemize}

\subsection{Maintenabilité}

Le projet doit rester :
\begin{itemize}
    \item modulaire ;
    \item lisible ;
    \item facilement évolutif ;
    \item réutilisable pour des évolutions futures.
\end{itemize}

\subsection{Disponibilité}

Le système doit assurer :
\begin{itemize}
    \item une disponibilité continue des fonctionnalités principales ;
    \item une séparation entre services critiques et traitements analytiques.
\end{itemize}

\section{Cas d’utilisation}

Les cas d’utilisation permettent de représenter les interactions principales entre les acteurs et le système.

\vspace{0.5cm}

\begin{figure}[H]
    \centering

    \maybeincludegraphics[width=0.9\textwidth]{figures/usecase-diagram.png}

    \fbox{
        \parbox[c][8cm][c]{0.85\textwidth}{
            \centering
            Placeholder :\\
            Diagramme de cas d’utilisation UML\\
            (Utilisateur + Administrateur + fonctionnalités principales)
        }
    }

    \caption{Diagramme de cas d’utilisation de SkillBridge}
\end{figure}

\subsection{Scénario : génération de recommandations}

Le scénario principal de la plateforme est la génération de recommandations pédagogiques.

\textbf{Acteur principal :} Utilisateur

\textbf{Précondition :}
L’utilisateur est authentifié.

\textbf{Scénario principal :}
\begin{enumerate}
    \item l’utilisateur crée une idée de projet ;
    \item le système analyse la description ;
    \item les compétences sont détectées ;
    \item les catégories pertinentes sont identifiées ;
    \item le moteur calcule les scores ;
    \item les cours sont classés ;
    \item les recommandations sont affichées.
\end{enumerate}

\textbf{Postcondition :}
Les résultats sont enregistrés dans l’historique des recommandations.

\section{Captures d’écran principales}

Cette section présente quelques interfaces principales de la plateforme.

\subsection{Page de connexion}

\begin{figure}[H]
    \centering

    \maybeincludegraphics[width=0.9\textwidth]{figures/login-page.png}

    \fbox{
        \parbox[c][7cm][c]{0.8\textwidth}{
            \centering
            Placeholder :\\
            Capture d’écran de la page de connexion
        }
    }

    \caption{Interface de connexion de SkillBridge}
\end{figure}

\subsection{Catalogue de cours}

\begin{figure}[H]
    \centering

    \maybeincludegraphics[width=0.95\textwidth]{figures/catalog-page.png}

    \fbox{
        \parbox[c][7cm][c]{0.85\textwidth}{
            \centering
            Placeholder :\\
            Capture du catalogue de cours\\
            avec filtres et pagination
        }
    }

    \caption{Catalogue des cours}
\end{figure}

\subsection{Recommandations de cours}

\begin{figure}[H]
    \centering

    \maybeincludegraphics[width=0.95\textwidth]{figures/recommendation-page.png}

    \fbox{
        \parbox[c][7cm][c]{0.85\textwidth}{
            \centering
            Placeholder :\\
            Capture des recommandations générées
        }
    }

    \caption{Résultats du moteur de recommandation}
\end{figure}

\subsection{Tableau de bord administrateur}

\begin{figure}[H]
    \centering

    \maybeincludegraphics[width=0.95\textwidth]{figures/admin-dashboard.png}

    \fbox{
        \parbox[c][7cm][c]{0.85\textwidth}{
            \centering
            Placeholder :\\
            Dashboard administrateur et analytics Big Data
        }
    }

    \caption{Tableau de bord analytique}
\end{figure}

\section{Conclusion}

Cette étude fonctionnelle a permis d’identifier les principaux besoins du système ainsi que les fonctionnalités essentielles de la plateforme SkillBridge.

L’analyse met en évidence la nécessité d’une architecture modulaire capable de séparer les traitements transactionnels et analytiques, tout en garantissant la sécurité, la maintenabilité et la scalabilité de l’application.

Le chapitre suivant présentera la conception et l’architecture globale du système.
% ---------------------------------------------------------------------------
\chapter{Conception et architecture du système}

\section{Introduction}

Après l’analyse des besoins fonctionnels et non fonctionnels, cette phase consiste à définir l’architecture globale de la plateforme SkillBridge ainsi que les choix techniques adoptés.


\section{Architecture globale}

L’objectif principal de cette architecture est de :
\begin{itemize}
    \item garantir une séparation claire des responsabilités ;
    \item assurer la sécurité des échanges ;
    \item faciliter la maintenabilité du système ;
    \item permettre l’intégration d’un pipeline Big Data indépendant ;
    \item assurer l’évolutivité de la plateforme.
\end{itemize}

L’architecture du projet repose sur une approche modulaire séparant :
\begin{itemize}
    \item le frontend React ;
    \item le backend Spring Boot ;
    \item la base de données PostgreSQL ;
    \item les composants Big Data Hadoop.
\end{itemize}

Le frontend React ne communique jamais directement avec :
\begin{itemize}
    \item PostgreSQL ;
    \item HDFS ;
    \item Hive ;
    \item HBase ;
    \item ou les conteneurs Docker.
\end{itemize}

Toutes les requêtes transitent par le backend Spring Boot qui joue le rôle de couche métier et de sécurité.

\vspace{0.5cm}

\begin{figure}[H]
    \centering

    \maybeincludegraphics[width=0.95\textwidth]{figures/global-architecture.png}

    \fbox{
        \parbox[c][9cm][c]{0.9\textwidth}{
            \centering
            Placeholder :\\
            Architecture globale SkillBridge\\
            \vspace{0.3cm}
            React → Spring Boot → PostgreSQL\\
            +
            Pipeline Hadoop (Flume, HDFS, Hive, MapReduce, HBase)
        }
    }

    \caption{Architecture globale de SkillBridge}
\end{figure}

\subsection{Séparation des couches}

Le système est organisé selon plusieurs couches indépendantes :

\begin{table}[H]
    \centering
    \caption{Organisation des couches du système}

    \begin{tabular}{|p{4cm}|p{9cm}|}
        \hline
        \textbf{Couche} & \textbf{Rôle} \\
        \hline

        Frontend &
        Interface utilisateur et interactions client. \\
        
        \hline

        Backend &
        API REST, logique métier, sécurité et accès aux données. \\
        
        \hline

        Base transactionnelle &
        Stockage des données applicatives temps réel. \\
        
        \hline

        Pipeline Big Data &
        Collecte, traitement et analyse des données massives. \\
        
        \hline

        Dashboard analytique &
        Visualisation des résultats et indicateurs Big Data. \\
        
        \hline
    \end{tabular}
\end{table}

\section{Architecture frontend}

Le frontend de SkillBridge a été développé avec React, Vite et TypeScript.

Cette couche assure :
\begin{itemize}
    \item l’affichage des interfaces utilisateur ;
    \item la gestion de navigation ;
    \item la communication avec l’API ;
    \item la gestion des sessions ;
    \item la protection des routes.
\end{itemize}

\subsection{Organisation du frontend}

Le frontend est structuré en plusieurs modules :

\begin{table}[H]
    \centering
    \caption{Organisation du frontend}

    \begin{tabular}{|p{4cm}|p{9cm}|}
        \hline
        \textbf{Dossier} & \textbf{Description} \\
        \hline

        src/pages &
        Pages principales de l’application. \\
        
        \hline

        src/components &
        Composants réutilisables de l’interface. \\
        
        \hline

        src/services &
        Communication avec l’API REST. \\
        
        \hline

        src/types &
        Définitions TypeScript et DTO frontend. \\
        
        \hline

        src/app &
        Gestion des routes et contexte global. \\
        
        \hline
    \end{tabular}
\end{table}

\subsection{Gestion des routes}

La navigation est assurée par React Router.

Le système utilise :
\begin{itemize}
    \item des routes publiques ;
    \item des routes protégées ;
    \item des routes réservées aux administrateurs.
\end{itemize}

\vspace{0.5cm}

\begin{figure}[H]
    \centering

    \maybeincludegraphics[width=0.85\textwidth]{figures/frontend-routing.png}

    \fbox{
        \parbox[c][6cm][c]{0.8\textwidth}{
            \centering
            Placeholder :\\
            Diagramme des routes frontend\\
            et mécanisme AuthGate/AdminGate
        }
    }

    \caption{Protection des routes frontend}
\end{figure}

\section{Architecture backend}

Le backend repose sur Spring Boot et suit une architecture orientée services.

Le backend centralise :
\begin{itemize}
    \item la logique métier ;
    \item la sécurité ;
    \item l’authentification ;
    \item la gestion des rôles ;
    \item l’accès aux données ;
    \item les endpoints REST ;
    \item l’intégration Big Data.
\end{itemize}

\subsection{Architecture des modules backend}

Le backend est organisé par domaines fonctionnels.

\begin{table}[H]
    \centering
    \caption{Modules backend principaux}

    \begin{tabular}{|p{4cm}|p{9cm}|}
        \hline
        \textbf{Module} & \textbf{Responsabilité} \\
        \hline

        security &
        JWT, Spring Security, OAuth et filtres. \\
        
        \hline

        user &
        Gestion des utilisateurs et rôles. \\
        
        \hline

        course &
        Gestion des cours et catalogue. \\
        
        \hline

        recommendation &
        Génération des recommandations. \\
        
        \hline

        projectidea &
        Gestion des projets utilisateur. \\
        
        \hline

        progress &
        Suivi pédagogique et progression. \\
        
        \hline

        bigdata &
        Intégration du pipeline analytique. \\
        
        \hline

        admin &
        Dashboard et statistiques administrateur. \\
        
        \hline
    \end{tabular}
\end{table}

\subsection{Architecture REST}

La communication entre frontend et backend repose sur une API REST utilisant le format JSON.

Les principales ressources exposées sont :
\begin{itemize}
    \item authentification ;
    \item utilisateurs ;
    \item catalogue ;
    \item projets ;
    \item recommandations ;
    \item progression ;
    \item analytics Big Data.
\end{itemize}

\subsection{Principaux endpoints REST}

\begin{table}[H]
\centering
\caption{Principaux endpoints REST de SkillBridge}
\begin{tabular}{|p{8.5cm}|p{2cm}|p{6cm}|}
\hline
\textbf{Endpoint} & \textbf{Accès} & \textbf{Rôle} \\
\hline
/api/auth/register & Public & Création de compte utilisateur \\
\hline
/api/auth/login & Public & Authentification JWT \\
\hline
/api/courses & Public & Consultation du catalogue \\
\hline
/api/projects & USER & Gestion des projets utilisateur \\
\hline
/api/projects/\{id\}/recommendations/generate & USER & Génération des recommandations \\
\hline
/api/saved-courses & USER & Sauvegarde des cours \\
\hline
/api/progress & USER & Suivi pédagogique \\
\hline
/api/admin/overview & ADMIN & Dashboard administrateur \\
\hline
/api/admin/bigdata/* & ADMIN & Supervision du pipeline Big Data \\
\hline
\end{tabular}
\end{table}

\vspace{0.5cm}

\begin{figure}[H]
    \centering

    \maybeincludegraphics[width=0.9\textwidth]{figures/api-architecture.png}

    \fbox{
        \parbox[c][6cm][c]{0.85\textwidth}{
            \centering
            Placeholder :\\
            Architecture API REST\\
            React $\leftrightarrow$ Spring Boot $\leftrightarrow$ PostgreSQL
        }
    }

    \caption{Architecture des échanges REST}
\end{figure}

\section{Sécurité de l’application}

La sécurité constitue un élément central de l’architecture de SkillBridge.

Le système adopte une approche backend-centric où toutes les décisions critiques sont prises côté serveur.

\subsection{Authentification}

Le système supporte :
\begin{itemize}
    \item l’authentification email/mot de passe ;
    \item l’authentification Google OAuth ;
    \item l’authentification GitHub OAuth.
\end{itemize}

Après authentification, le backend génère un JWT utilisé pour sécuriser les requêtes suivantes.

\subsection{Gestion des rôles}

Deux rôles principaux sont définis :
\begin{itemize}
    \item USER ;
    \item ADMIN.
\end{itemize}

Les autorisations sont contrôlées par Spring Security.

\subsection{Mesures de sécurité}

Plusieurs mécanismes ont été mis en place :
\begin{itemize}
    \item hashage des mots de passe avec BCrypt ;
    \item protection des endpoints ;
    \item validation des rôles ;
    \item gestion CORS ;
    \item limitation des tentatives de connexion ;
    \item protection contre les accès non autorisés.
\end{itemize}

\vspace{0.5cm}

\begin{figure}[H]
    \centering

    \maybeincludegraphics[width=0.9\textwidth]{figures/jwt-flow.png}

    \fbox{
        \parbox[c][7cm][c]{0.85\textwidth}{
            \centering
            Placeholder :\\
            Schéma du flux JWT\\
            Login → Token → Authorization Header
        }
    }

    \caption{Flux d’authentification JWT}
\end{figure}

\section{Modèle de données}

Le stockage principal de l’application repose sur PostgreSQL via Supabase.

Le modèle de données est centré autour :
\begin{itemize}
    \item des utilisateurs ;
    \item des projets ;
    \item des cours ;
    \item des compétences ;
    \item des recommandations.
\end{itemize}

\subsection{Tables principales}

\begin{table}[H]
    \centering
    \caption{Principales tables de la base de données}

    \begin{tabular}{|p{4cm}|p{9cm}|}
        \hline
        \textbf{Table} & \textbf{Description} \\
        \hline

        users &
        Comptes utilisateurs et rôles. \\
        
        \hline

        courses &
        Catalogue des cours disponibles. \\
        
        \hline

        skills &
        Compétences détectées et associées. \\
        
        \hline

        project\_ideas &
        Projets créés par les utilisateurs. \\
        
        \hline

        recommendation\_results &
        Résultats générés par le moteur. \\
        
        \hline

        saved\_courses &
        Cours sauvegardés par utilisateur. \\
        
        \hline

        course\_progress &
        Progression pédagogique des utilisateurs. \\
        
        \hline
    \end{tabular}
\end{table}

\subsection{Relations entre entités}

Le système utilise plusieurs relations :
\begin{itemize}
    \item utilisateur → projets ;
    \item projets → recommandations ;
    \item cours $\leftrightarrow$ compétences ;
    \item utilisateur → progression ;
    \item utilisateur → favoris.
\end{itemize}

\vspace{0.5cm}

\begin{figure}[H]
    \centering

    \maybeincludegraphics[width=0.95\textwidth]{figures/database-model.png}

    \fbox{
        \parbox[c][8cm][c]{0.9\textwidth}{
            \centering
            Placeholder :\\
            Diagramme relationnel / MCD\\
            de la base de données SkillBridge
        }
    }

    \caption{Modèle relationnel simplifié}
\end{figure}

\section{Technologies utilisées}

Le projet repose sur plusieurs technologies complémentaires.

\begin{table}[H]
    \centering
    \caption{Technologies principales utilisées}

    \begin{tabular}{|p{4cm}|p{4cm}|p{5cm}|}
        \hline
        \textbf{Couche} & \textbf{Technologie} & \textbf{Rôle} \\
        \hline

        Frontend &
        React + Vite &
        Interface utilisateur dynamique. \\
        
        \hline

        Backend &
        Spring Boot &
        API REST et logique métier. \\
        
        \hline

        Sécurité &
        Spring Security + JWT &
        Authentification et autorisation. \\
        
        \hline

        Base de données &
        PostgreSQL / Supabase &
        Stockage transactionnel. \\
        
        \hline

        Big Data &
        Hadoop Ecosystem &
        Traitement analytique distribué. \\
        
        \hline

        Streaming &
        Flume &
        Collecte des événements. \\
        
        \hline

        SQL analytique &
        Hive &
        Analyse SQL sur HDFS. \\
        
        \hline

        Batch &
        MapReduce &
        Traitement distribué. \\
        
        \hline

        NoSQL &
        HBase &
        Stockage analytique rapide. \\
        
        \hline
    \end{tabular}
\end{table}

\section{Conclusion}

Ce chapitre a présenté la conception générale et l’architecture de la plateforme SkillBridge.

L’architecture proposée repose sur une séparation claire entre :
\begin{itemize}
    \item la couche applicative temps réel ;
    \item la couche de sécurité ;
    \item la couche de persistance ;
    \item et la couche analytique Big Data.
\end{itemize}

Cette organisation permet de garantir la modularité, la maintenabilité et l’évolutivité du système.

Le chapitre suivant présentera le moteur de recommandation ainsi que les mécanismes de détection et de scoring utilisés par la plateforme.
% ---------------------------------------------------------------------------
\chapter{Moteur de recommandation}

\section{Introduction}

Le moteur de recommandation constitue l’un des composants centraux de la plateforme SkillBridge. Son objectif est de transformer une idée de projet exprimée en langage naturel en une liste ordonnée de cours pertinents.

Contrairement aux systèmes reposant exclusivement sur des modèles d’intelligence artificielle opaques, le moteur développé dans SkillBridge adopte une approche déterministe et explicable. Chaque recommandation produite peut être justifiée à partir :
\begin{itemize}
    \item des mots-clés détectés ;
    \item des compétences identifiées ;
    \item des catégories associées ;
    \item et du score calculé.
\end{itemize}

Cette approche permet :
\begin{itemize}
    \item d’améliorer la transparence du système ;
    \item de faciliter l’interprétation des résultats ;
    \item de rendre les recommandations auditables ;
    \item et de simplifier l’analyse académique du projet.
\end{itemize}

\section{Principe général du moteur}

Le fonctionnement général du moteur de recommandation suit plusieurs étapes successives.

\begin{enumerate}
    \item L’utilisateur saisit une idée de projet.
    
    \item Le texte est normalisé et analysé.
    
    \item Les mots-clés importants sont extraits.
    
    \item Les compétences correspondantes sont détectées.
    
    \item Les catégories liées au projet sont identifiées.
    
    \item Les cours candidats sont sélectionnés.
    
    \item Un score est calculé pour chaque cours.
    
    \item Les résultats sont classés et retournés à l’utilisateur.
\end{enumerate}

\vspace{0.5cm}

\begin{figure}[H]
    \centering

    \maybeincludegraphics[width=0.95\textwidth]{figures/recommendation-pipeline.png}

    \fbox{
        \parbox[c][8cm][c]{0.9\textwidth}{
            \centering
            Placeholder :\\
            Pipeline du moteur de recommandation\\
            \vspace{0.3cm}
            Project Idea → Keyword Extraction → Skill Detection\\
            → Category Matching → Scoring → Ranked Courses
        }
    }

    \caption{Pipeline du moteur de recommandation}
\end{figure}

\section{Prétraitement et normalisation du texte}

Avant de générer des recommandations, le texte fourni par l’utilisateur doit être normalisé afin de faciliter les traitements.

Les principales opérations réalisées sont :
\begin{itemize}
    \item conversion en minuscules ;
    \item suppression des caractères spéciaux ;
    \item nettoyage des espaces ;
    \item tokenisation des mots ;
    \item filtrage des mots peu significatifs.
\end{itemize}

Cette étape permet de réduire les variations syntaxiques et d’améliorer la détection des compétences.

\section{Détection des compétences}

Après normalisation, le système tente d’identifier les compétences présentes dans la description du projet.

Le moteur repose sur :
\begin{itemize}
    \item des mots-clés ;
    \item des correspondances textuelles ;
    \item des catégories techniques ;
    \item un catalogue enrichi de compétences.
\end{itemize}

Par exemple :
\begin{itemize}
    \item le mot \textit{React} peut être associé au développement frontend ;
    \item le mot \textit{Java} peut être associé au backend ;
    \item le mot \textit{Hadoop} peut être associé au Big Data.
\end{itemize}

Les compétences détectées sont ensuite comparées aux compétences associées aux cours du catalogue.

\vspace{0.5cm}

\begin{figure}[H]
    \centering

    \maybeincludegraphics[width=0.85\textwidth]{figures/skill-detection.png}

    \fbox{
        \parbox[c][7cm][c]{0.8\textwidth}{
            \centering
            Placeholder :\\
            Schéma de détection des compétences\\
            depuis une idée de projet
        }
    }

    \caption{Détection des compétences}
\end{figure}

\section{Détection des catégories}

En plus des compétences, le système identifie des catégories générales correspondant au projet utilisateur.

Les principales catégories utilisées sont :
\begin{itemize}
    \item Web Development ;
    \item Backend Development ;
    \item Big Data ;
    \item DevOps ;
    \item Machine Learning ;
    \item Databases ;
    \item Application Security.
\end{itemize}

La détection des catégories permet :
\begin{itemize}
    \item d’améliorer la pertinence des recommandations ;
    \item de filtrer les cours ;
    \item d’enrichir le calcul du score final.
\end{itemize}

\section{Algorithme de scoring}

Le classement des cours repose sur un système de scoring déterministe.

Chaque cours reçoit un score calculé à partir de plusieurs signaux :
\begin{itemize}
    \item correspondance du titre ;
    \item correspondance des compétences ;
    \item correspondance des catégories ;
    \item bonus de popularité.
\end{itemize}

\subsection{Composition du score}

\begin{table}[H]
    \centering
    \caption{Composition du score de recommandation}

    \begin{tabular}{|p{4cm}|p{3cm}|p{7cm}|}
        \hline
        \textbf{Signal} & \textbf{Score max} & \textbf{Description} \\
        \hline

        Titre &
        30 &
        Présence des mots-clés du projet dans le titre du cours. \\
        
        \hline

        Compétences &
        40 &
        Correspondance entre les compétences détectées et les compétences du cours. \\
        
        \hline

        Catégorie &
        20 &
        Correspondance entre la catégorie détectée et celle du cours. \\
        
        \hline

        Bonus &
        10 &
        Popularité et indicateurs analytiques disponibles. \\
        
        \hline
    \end{tabular}
\end{table}

\subsection{Formule générale}

Le score total est calculé selon la formule suivante :

\[
Score_{total} =
Score_{titre}
+
Score_{competences}
+
Score_{categorie}
+
Bonus
\]

Le score final est limité à une valeur maximale de 100.

\subsection{Exemple de calcul}

Considérons le projet suivant :

\begin{quote}
\textit{
"I want to create a complete website using React and Java"
}
\end{quote}

Le système peut produire le résultat suivant :

\begin{table}[H]
    \centering
    \caption{Exemple de calcul de score}

    \begin{tabular}{|p{5cm}|p{3cm}|}
        \hline
        \textbf{Composant} & \textbf{Valeur} \\
        \hline

        Title Match Score &
        24 \\
        
        \hline

        Skill Match Score &
        36 \\
        
        \hline

        Category Match Score &
        20 \\
        
        \hline

        Popularity Bonus &
        6.4 \\
        
        \hline

        Total &
        86.4 \\
        
        \hline
    \end{tabular}
\end{table}

\section{Classement des recommandations}

Une fois les scores calculés :
\begin{itemize}
    \item les cours sont triés ;
    \item les meilleurs résultats sont sélectionnés ;
    \item les explications sont générées ;
    \item les résultats sont sauvegardés.
\end{itemize}

Le système peut également limiter le nombre de recommandations retournées.

\section{Explicabilité des recommandations}

L’un des objectifs principaux du moteur est de rendre les recommandations compréhensibles.

Pour chaque résultat, le système conserve :
\begin{itemize}
    \item le score total ;
    \item le détail du score ;
    \item les compétences détectées ;
    \item les catégories associées ;
    \item les raisons de la recommandation.
\end{itemize}

Cette approche permet :
\begin{itemize}
    \item d’améliorer la confiance utilisateur ;
    \item de simplifier le débogage ;
    \item de rendre les résultats auditables ;
    \item de justifier les recommandations devant un jury académique.
\end{itemize}

\subsection{Snapshots de recommandations}

Chaque génération crée :
\begin{itemize}
    \item un \texttt{RecommendationSnapshot} ;
    \item plusieurs \texttt{RecommendationResult}.
\end{itemize}

Les snapshots permettent :
\begin{itemize}
    \item de conserver l’historique ;
    \item d’analyser les recommandations générées ;
    \item de reproduire les résultats ultérieurement.
\end{itemize}

\vspace{0.5cm}

\begin{figure}[H]
    \centering

    \maybeincludegraphics[width=0.9\textwidth]{figures/recommendation-snapshot.png}

    \fbox{
        \parbox[c][6cm][c]{0.85\textwidth}{
            \centering
            Placeholder :\\
            Schéma snapshots et recommendation results
        }
    }

    \caption{Historisation des recommandations}
\end{figure}

\section{Intégration avec le pipeline Big Data}

Le moteur de recommandation est également connecté au pipeline analytique.

Lorsqu’une recommandation est générée :
\begin{itemize}
    \item un événement est écrit dans \texttt{events.log} ;
    \item Flume collecte cet événement ;
    \item les données sont envoyées vers HDFS ;
    \item Hive et MapReduce peuvent analyser les recommandations générées.
\end{itemize}

Cette intégration permet :
\begin{itemize}
    \item l’analyse des comportements utilisateurs ;
    \item l’identification des technologies populaires ;
    \item la génération de statistiques analytiques ;
    \item l’amélioration future du système.
\end{itemize}

\section{Limites du moteur actuel}

Malgré son efficacité, le moteur actuel présente certaines limites.

\subsection{Approche déterministe}

Le moteur repose principalement sur des règles fixes :
\begin{itemize}
    \item matching textuel ;
    \item règles de scoring ;
    \item catégories prédéfinies.
\end{itemize}

Cette approche peut limiter :
\begin{itemize}
    \item la compréhension sémantique ;
    \item la flexibilité ;
    \item la personnalisation avancée.
\end{itemize}

\subsection{Absence de Machine Learning}

Le système n’utilise pas encore :
\begin{itemize}
    \item de collaborative filtering ;
    \item d’embeddings ;
    \item de réseaux neuronaux ;
    \item de modèles NLP avancés ;
    \item de modèles LLM.
\end{itemize}

\section{Perspectives d’amélioration}

Plusieurs améliorations peuvent être envisagées pour les futures versions du projet.

\subsection{Intelligence artificielle}

Le moteur pourrait intégrer :
\begin{itemize}
    \item des modèles NLP ;
    \item des embeddings sémantiques ;
    \item des modèles Transformers ;
    \item des LLM ;
    \item des systèmes hybrides ML + règles.
\end{itemize}

\subsection{Personnalisation avancée}

Le système pourrait également exploiter :
\begin{itemize}
    \item l’historique utilisateur ;
    \item les interactions précédentes ;
    \item les préférences pédagogiques ;
    \item les statistiques comportementales.
\end{itemize}

\subsection{Analyse prédictive}

Les données analytiques collectées pourraient servir à :
\begin{itemize}
    \item prédire les besoins utilisateurs ;
    \item identifier les tendances ;
    \item améliorer automatiquement le classement des cours.
\end{itemize}

\section{Conclusion}

Ce chapitre a présenté le moteur de recommandation développé dans SkillBridge.

Le système repose sur une approche déterministe et explicable permettant :
\begin{itemize}
    \item la détection des compétences ;
    \item le calcul de scores pertinents ;
    \item le classement des cours ;
    \item la traçabilité des recommandations.
\end{itemize}

L’intégration avec le pipeline Big Data permet également de collecter et analyser les comportements utilisateurs afin d’enrichir les fonctionnalités analytiques de la plateforme.

Le chapitre suivant présentera l’architecture Big Data ainsi que les différents composants Hadoop intégrés dans le projet.

% ---------------------------------------------------------------------------
\chapter{Pipeline Big Data et analytique}

\section{Introduction}

En plus de la partie applicative full-stack, le projet SkillBridge intègre un pipeline Big Data basé sur l’écosystème Hadoop.

L’objectif principal de cette couche analytique est de :
\begin{itemize}
    \item collecter les événements générés par les utilisateurs ;
    \item traiter des données volumineuses ;
    \item stocker les données dans un environnement distribué ;
    \item exécuter des traitements analytiques ;
    \item produire des statistiques exploitables ;
    \item alimenter le tableau de bord administrateur.
\end{itemize}

Le pipeline Big Data a été conçu de manière indépendante de l’application transactionnelle afin de préserver les performances de l’API temps réel.

Le principe fondamental adopté dans le projet est le suivant :
\begin{itemize}
    \item PostgreSQL/Supabase gère les opérations transactionnelles temps réel ;
    \item Hadoop gère les traitements analytiques batch et streaming.
\end{itemize}

Cette séparation permet de construire une architecture scalable et adaptée aux traitements distribués.

\section{Objectif du pipeline Big Data}

Le pipeline Big Data permet de démontrer plusieurs concepts importants liés au traitement distribué des données.

Les objectifs principaux sont :
\begin{itemize}
    \item ingestion de données externes ;
    \item nettoyage et enrichissement du catalogue ;
    \item collecte des événements utilisateurs ;
    \item stockage distribué avec HDFS ;
    \item analyse SQL avec Hive ;
    \item traitement batch avec MapReduce ;
    \item stockage NoSQL avec HBase ;
    \item visualisation analytique dans le dashboard administrateur.
\end{itemize}

Le pipeline est également utilisé pour enrichir les fonctionnalités du moteur de recommandation.

\section{Architecture générale du pipeline}

L’architecture Big Data repose sur deux flux principaux :
\begin{itemize}
    \item un flux batch ;
    \item un flux streaming temps réel.
\end{itemize}

\vspace{0.5cm}

\begin{figure}[H]
    \centering

    \maybeincludegraphics[width=0.95\textwidth]{figures/bigdata-pipeline.png}

    \fbox{
        \parbox[c][10cm][c]{0.9\textwidth}{
            \centering
            Placeholder :\\
            Architecture complète du pipeline Big Data\\
            \vspace{0.4cm}
            ZIP/CSV → Python → PostgreSQL Mirror → Sqoop → HDFS\\
            \vspace{0.2cm}
            React → Spring Boot → events.log → Flume → HDFS\\
            \vspace{0.2cm}
            HDFS → Hive / MapReduce / HBase → Dashboard Admin
        }
    }

    \caption{Architecture du pipeline Big Data}
\end{figure}

\section{Flux batch de préparation du catalogue}

Le flux batch permet de construire un catalogue de cours propre et exploitable.

Les données initiales proviennent de plusieurs datasets externes au format :
\begin{itemize}
    \item ZIP ;
    \item CSV ;
    \item JSON.
\end{itemize}

\subsection{Fusion et nettoyage des données}

Des scripts Python ont été utilisés pour :
\begin{itemize}
    \item fusionner les datasets ;
    \item supprimer les doublons ;
    \item nettoyer les données ;
    \item normaliser les catégories ;
    \item détecter les compétences ;
    \item générer un catalogue unifié.
\end{itemize}

Le résultat final contient :
\begin{itemize}
    \item 17072 cours ;
    \item 13647 compétences ;
    \item 459 fournisseurs ;
    \item 12 catégories.
\end{itemize}

\subsection{Statistiques du catalogue}

\begin{table}[H]
    \centering
    \caption{Statistiques du catalogue enrichi}

    \begin{tabular}{|p{6cm}|p{4cm}|}
        \hline
        \textbf{Élément} & \textbf{Valeur} \\
        \hline

        Nombre total de cours &
        17072 \\
        
        \hline

        Nombre de compétences &
        13647 \\
        
        \hline

        Nombre de catégories &
        12 \\
        
        \hline

        Nombre de fournisseurs &
        459 \\
        
        \hline

        Relations course\_skills &
        104006 \\
        
        \hline
    \end{tabular}
\end{table}

\subsection{Distribution des catégories}

Les catégories les plus représentées dans le catalogue sont :
\begin{itemize}
    \item Machine Learning ;
    \item Backend Development ;
    \item Product and UX ;
    \item Application Security ;
    \item Cloud Computing ;
    \item Big Data ;
    \item DevOps.
\end{itemize}

\vspace{0.5cm}

\begin{figure}[H]
    \centering

    \maybeincludegraphics[width=0.85\textwidth]{figures/catalog-stats.png}

    \fbox{
        \parbox[c][7cm][c]{0.8\textwidth}{
            \centering
            Placeholder :\\
            Graphiques statistiques du catalogue\\
            (catégories, niveaux, compétences)
        }
    }

    \caption{Statistiques du catalogue enrichi}
\end{figure}

\section{PostgreSQL mirror local}

Le pipeline analytique n’utilise pas directement la base PostgreSQL officielle de l’application.

Une copie locale appelée \textbf{PostgreSQL mirror} est utilisée afin de :
\begin{itemize}
    \item éviter d’impacter la base transactionnelle ;
    \item fournir une source stable pour Sqoop ;
    \item faciliter les tests Hadoop ;
    \item permettre la reconstruction du pipeline sans risque.
\end{itemize}

Le mirror contient :
\begin{itemize}
    \item les cours ;
    \item les compétences ;
    \item les catégories ;
    \item les projets ;
    \item les progressions ;
    \item les données analytiques utiles.
\end{itemize}

\section{Stockage distribué avec HDFS}

Le stockage distribué repose sur HDFS (Hadoop Distributed File System).

HDFS permet :
\begin{itemize}
    \item le stockage distribué des fichiers ;
    \item la réplication des blocs ;
    \item la tolérance aux pannes ;
    \item la scalabilité horizontale.
\end{itemize}

\subsection{NameNode et DataNodes}

Le cluster Hadoop local est composé :
\begin{itemize}
    \item d’un NameNode ;
    \item de plusieurs DataNodes.
\end{itemize}

Le NameNode gère :
\begin{itemize}
    \item les métadonnées ;
    \item les chemins HDFS ;
    \item la localisation des blocs.
\end{itemize}

Les DataNodes assurent :
\begin{itemize}
    \item le stockage réel des données ;
    \item la réplication des blocs.
\end{itemize}

Le projet utilise un cluster avec deux DataNodes afin de démontrer la scalabilité horizontale.

\vspace{0.5cm}

\begin{figure}[H]
    \centering

    \maybeincludegraphics[width=0.9\textwidth]{figures/hdfs-cluster.png}

    \fbox{
        \parbox[c][7cm][c]{0.85\textwidth}{
            \centering
            Placeholder :\\
            Architecture HDFS\\
            NameNode + 2 DataNodes
        }
    }

    \caption{Architecture du cluster HDFS}
\end{figure}

\section{Importation batch avec Sqoop}

Sqoop est utilisé pour importer les données relationnelles du PostgreSQL mirror vers HDFS.

Le flux suivi est le suivant :

\[
PostgreSQL \rightarrow Sqoop \rightarrow HDFS
\]

Sqoop permet :
\begin{itemize}
    \item d’automatiser les imports relationnels ;
    \item de transférer les tables PostgreSQL vers Hadoop ;
    \item de générer les fichiers exploitables dans HDFS.
\end{itemize}

\subsection{Tables importées}

Les principales tables importées sont :
\begin{itemize}
    \item courses ;
    \item skills ;
    \item categories ;
    \item providers ;
    \item project\_ideas ;
    \item saved\_courses ;
    \item course\_progress.
\end{itemize}

\subsection{Résultats d’import}

Chaque import réussi génère :
\begin{itemize}
    \item un fichier \texttt{part-m-00000} ;
    \item un fichier \texttt{\_SUCCESS}.
\end{itemize}

La présence du fichier \texttt{\_SUCCESS} confirme la réussite de l’import batch.

\vspace{0.5cm}

\begin{figure}[H]
    \centering

    \maybeincludegraphics[width=0.95\textwidth]{figures/sqoop-import.png}

    \fbox{
        \parbox[c][7cm][c]{0.9\textwidth}{
            \centering
            Placeholder :\\
            Capture HDFS après import Sqoop\\
            avec fichiers \_SUCCESS
        }
    }

    \caption{Import batch avec Sqoop}
\end{figure}

\section{Collecte streaming avec Flume}

Le pipeline utilise Flume pour collecter les événements générés par l’application web.

\subsection{Principe de fonctionnement}

Lorsqu’un utilisateur effectue une action :
\begin{itemize}
    \item Spring Boot écrit un événement JSON dans \texttt{events.log} ;
    \item Flume surveille ce fichier ;
    \item les événements sont envoyés vers HDFS.
\end{itemize}

Les événements collectés incluent :
\begin{itemize}
    \item COURSE\_SEARCH ;
    \item COURSE\_CLICK ;
    \item COURSE\_SAVE ;
    \item PROJECT\_CREATED ;
    \item PROJECT\_RECOMMENDATION.
\end{itemize}

\subsection{Architecture Flume}

Le pipeline Flume repose sur :
\begin{itemize}
    \item une source ;
    \item un channel mémoire ;
    \item un sink HDFS.
\end{itemize}

\vspace{0.5cm}

\begin{figure}[H]
    \centering

    \maybeincludegraphics[width=0.9\textwidth]{figures/flume-pipeline.png}

    \fbox{
        \parbox[c][7cm][c]{0.85\textwidth}{
            \centering
            Placeholder :\\
            Architecture Flume\\
            events.log → Flume → HDFS
        }
    }

    \caption{Collecte streaming avec Flume}
\end{figure}

\section{Analyse SQL avec Hive}

Hive permet d’exécuter des requêtes SQL sur les fichiers stockés dans HDFS.

Les principales tables Hive utilisées sont :
\begin{itemize}
    \item hive\_courses ;
    \item hive\_events ;
    \item hive\_saved\_courses ;
    \item hive\_course\_progress.
\end{itemize}

\subsection{Analyse des événements}

Les logs collectés par Flume sont stockés sous forme JSON brute.

Hive permet d’extraire les champs JSON avec :
\begin{itemize}
    \item \texttt{get\_json\_object()} ;
    \item des requêtes analytiques ;
    \item des agrégations statistiques.
\end{itemize}

\subsection{Résultats analytiques}

Les analyses Hive permettent par exemple :
\begin{itemize}
    \item de compter les événements ;
    \item d’identifier les recherches populaires ;
    \item d’analyser les comportements utilisateurs.
\end{itemize}

\vspace{0.5cm}

\begin{figure}[H]
    \centering

    \maybeincludegraphics[width=0.9\textwidth]{figures/hive-analysis.png}

    \fbox{
        \parbox[c][7cm][c]{0.85\textwidth}{
            \centering
            Placeholder :\\
            Requêtes Hive et résultats analytiques
        }
    }

    \caption{Analyse SQL avec Hive}
\end{figure}

\section{Traitement distribué avec MapReduce}

Le projet utilise un job MapReduce appelé :
\texttt{TopSearchKeywordsJob}.

Ce traitement analyse les événements collectés afin d’identifier les mots-clés les plus fréquents.

\subsection{Fonctionnement}

Le traitement suit les étapes classiques :
\begin{itemize}
    \item lecture des événements ;
    \item extraction des mots ;
    \item phase Map ;
    \item phase Shuffle ;
    \item phase Reduce ;
    \item génération des résultats.
\end{itemize}

\subsection{Résultats obtenus}

Les principaux mots-clés détectés incluent :
\begin{itemize}
    \item react ;
    \item python ;
    \item machine ;
    \item spring ;
    \item frontend ;
    \item hadoop ;
    \item hive.
\end{itemize}

\vspace{0.5cm}

\begin{figure}[H]
    \centering

    \maybeincludegraphics[width=0.9\textwidth]{figures/mapreduce-flow.png}

    \fbox{
        \parbox[c][7cm][c]{0.85\textwidth}{
            \centering
            Placeholder :\\
            Schéma MapReduce\\
            Map → Shuffle → Reduce
        }
    }

    \caption{Traitement distribué MapReduce}
\end{figure}

\section{Stockage analytique avec HBase}

HBase est utilisé comme couche NoSQL du pipeline.

La table principale utilisée est :
\texttt{course\_stats}.

\subsection{Informations stockées}

Les statistiques stockées incluent :
\begin{itemize}
    \item nombre de clics ;
    \item nombre de sauvegardes ;
    \item progression moyenne ;
    \item titre du cours.
\end{itemize}

\subsection{Structure des colonnes}

Les familles de colonnes utilisées sont :
\begin{itemize}
    \item meta:title ;
    \item activity:clicks ;
    \item activity:saves ;
    \item activity:avg\_progress.
\end{itemize}

HBase permet une lecture rapide basée sur la clé \texttt{courseId}.

\vspace{0.5cm}

\begin{figure}[H]
    \centering

    \maybeincludegraphics[width=0.9\textwidth]{figures/hbase-scan.png}

    \fbox{
        \parbox[c][7cm][c]{0.85\textwidth}{
            \centering
            Placeholder :\\
            Résultat scan HBase\\
            course\_stats
        }
    }

    \caption{Stockage analytique avec HBase}
\end{figure}

\section{Dashboard analytique}

Les résultats du pipeline sont exposés via :
\begin{itemize}
    \item des fichiers JSON ;
    \item des endpoints Spring Boot ;
    \item un dashboard administrateur.
\end{itemize}

Le dashboard permet :
\begin{itemize}
    \item de visualiser l’état du pipeline ;
    \item d’afficher les statistiques analytiques ;
    \item de consulter les événements ;
    \item de vérifier les traitements Hadoop.
\end{itemize}

\section{Intégration avec l’application web}

Le frontend React ne communique jamais directement avec Hadoop.

Le fonctionnement adopté est :
\[
React \rightarrow Spring Boot \rightarrow BigDataOutputs
\]

Spring Boot :
\begin{itemize}
    \item lit les fichiers générés ;
    \item transforme les résultats ;
    \item expose les données via des endpoints REST.
\end{itemize}

Cette approche garantit :
\begin{itemize}
    \item la sécurité ;
    \item la modularité ;
    \item la stabilité de l’architecture.
\end{itemize}

\section{Limites actuelles}

Malgré son bon fonctionnement, le pipeline présente certaines limites.

\begin{itemize}
    \item Le pipeline reste terminal-first et non totalement automatisé.
    
    \item HBase peut nécessiter un redémarrage avant démonstration.
    
    \item Les événements Hive sont stockés sous forme JSON brute.
    
    \item Certains traitements restent semi-manuels.
    
    \item Le système n’utilise pas encore Spark ou Kafka.
\end{itemize}

\section{Perspectives d’évolution}

Plusieurs améliorations peuvent être envisagées :

\begin{itemize}
    \item automatisation complète du pipeline ;
    
    \item intégration de Spark Streaming ;
    
    \item utilisation de Kafka ;
    
    \item analyse temps réel avancée ;
    
    \item intégration IA et Machine Learning ;
    
    \item dashboards analytiques plus avancés.
\end{itemize}

\section{Conclusion}

Ce chapitre a présenté l’architecture et le fonctionnement du pipeline Big Data intégré dans SkillBridge.

Le pipeline combine plusieurs technologies complémentaires :
\begin{itemize}
    \item Python ;
    \item PostgreSQL mirror ;
    \item Sqoop ;
    \item HDFS ;
    \item Flume ;
    \item Hive ;
    \item MapReduce ;
    \item HBase.
\end{itemize}

Cette architecture permet de démontrer :
\begin{itemize}
    \item l’ingestion de données ;
    \item le stockage distribué ;
    \item le traitement analytique ;
    \item la collecte streaming ;
    \item et la visualisation des résultats.
\end{itemize}

Le chapitre suivant présentera le déploiement du système, les validations réalisées ainsi que les résultats obtenus.
% ---------------------------------------------------------------------------
\chapter{Déploiement, validation et résultats}

\section{Introduction}

Après le développement de la plateforme et l’intégration du pipeline Big Data, plusieurs phases de validation ont été réalisées afin de vérifier :
\begin{itemize}
    \item le bon fonctionnement de l’application ;
    \item la stabilité des services ;
    \item la communication entre les composants ;
    \item le fonctionnement du pipeline analytique ;
    \item la cohérence des résultats générés.
\end{itemize}

Ce chapitre présente l’environnement de développement utilisé, les étapes de déploiement, les validations fonctionnelles ainsi que les résultats obtenus.

\section{Environnement de développement}

Le projet a été développé dans un environnement local basé sur Docker et plusieurs outils de développement modernes.

\subsection{Environnement logiciel}

\begin{table}[H]
    \centering
    \caption{Environnement de développement}

    \begin{tabular}{|p{5cm}|p{6cm}|}
        \hline
        \textbf{Composant} & \textbf{Technologie utilisée} \\
        \hline

        Système d’exploitation &
        Windows \\
        
        \hline

        Frontend &
        React + Vite + TypeScript \\
        
        \hline

        Backend &
        Spring Boot + Java 21 \\
        
        \hline

        Base de données &
        PostgreSQL / Supabase \\
        
        \hline

        Big Data &
        Hadoop Ecosystem \\
        
        \hline

        Conteneurisation &
        Docker + Docker Compose \\
        
        \hline

        Scripts &
        Python + PowerShell \\
        
        \hline
    \end{tabular}
\end{table}

\subsection{Architecture de déploiement}

Le projet est séparé en plusieurs modules indépendants :
\begin{itemize}
    \item frontend React ;
    \item backend Spring Boot ;
    \item PostgreSQL ;
    \item cluster Hadoop ;
    \item pipeline analytique.
\end{itemize}

\vspace{0.5cm}

\begin{figure}[H]
    \centering

    \maybeincludegraphics[width=0.95\textwidth]{figures/deployment-architecture.png}

    \fbox{
        \parbox[c][8cm][c]{0.9\textwidth}{
            \centering
            Placeholder :\\
            Architecture de déploiement\\
            Frontend + Backend + Docker Hadoop Stack
        }
    }

    \caption{Architecture de déploiement du projet}
\end{figure}

\section{Déploiement de l’application}

\subsection{Déploiement du backend}

Le backend Spring Boot expose l’ensemble des endpoints REST de l’application.

Le lancement du backend permet :
\begin{itemize}
    \item l’accès à l’API ;
    \item l’authentification ;
    \item la communication avec PostgreSQL ;
    \item l’exposition des endpoints Big Data.
\end{itemize}

\vspace{0.3cm}

\textbf{Placeholder :}
\begin{itemize}
    \item capture du backend démarré ;
    \item capture terminal Spring Boot ;
    \item logs de démarrage.
\end{itemize}

\subsection{Déploiement du frontend}

Le frontend React permet l’accès utilisateur à la plateforme.

Le frontend fournit :
\begin{itemize}
    \item les interfaces utilisateur ;
    \item le catalogue de cours ;
    \item les projets ;
    \item les recommandations ;
    \item le dashboard administrateur.
\end{itemize}

\vspace{0.3cm}

\textbf{Placeholder :}
\begin{itemize}
    \item capture du frontend ;
    \item interface login ;
    \item interface dashboard.
\end{itemize}

\subsection{Déploiement de la stack Big Data}

Le pipeline Big Data est exécuté dans des conteneurs Docker.

Les principaux services démarrés sont :
\begin{itemize}
    \item NameNode ;
    \item DataNodes ;
    \item Hive ;
    \item Flume ;
    \item HBase ;
    \item PostgreSQL mirror ;
    \item Sqoop client.
\end{itemize}

\vspace{0.3cm}

\textbf{Placeholder :}
\begin{itemize}
    \item capture docker compose ps ;
    \item capture conteneurs actifs ;
    \item capture cluster Hadoop.
\end{itemize}

\section{Validation fonctionnelle}

\subsection{Tests automatisés}

Le backend contient plusieurs tests automatisés permettant de vérifier le bon fonctionnement des composants critiques.

\begin{table}[H]
\centering
\caption{Tests automatisés principaux}
\begin{tabular}{|p{6cm}|p{7cm}|}
\hline
\textbf{Test} & \textbf{Objectif} \\
\hline
LoginAttemptServiceTests & Vérification du blocage après plusieurs tentatives échouées \\
\hline
RecommendationServiceTests & Validation de la normalisation et extraction des mots-clés \\
\hline
SkillBridgeApplicationTests & Vérification du chargement du contexte Spring Boot \\
\hline
\end{tabular}
\end{table}

Cette phase vise à vérifier le bon fonctionnement des fonctionnalités principales de l’application.

\subsection{Authentification}

Les tests réalisés ont permis de valider :
\begin{itemize}
    \item l’inscription utilisateur ;
    \item la connexion ;
    \item la génération des JWT ;
    \item la protection des routes ;
    \item la gestion des rôles.
\end{itemize}

\vspace{0.3cm}

\textbf{Placeholder :}
\begin{itemize}
    \item capture login réussi ;
    \item capture JWT ;
    \item capture accès admin.
\end{itemize}

\subsection{Gestion du catalogue}

Les fonctionnalités validées incluent :
\begin{itemize}
    \item affichage des cours ;
    \item pagination ;
    \item filtres ;
    \item recherche ;
    \item consultation des détails.
\end{itemize}

\vspace{0.3cm}

\textbf{Placeholder :}
\begin{itemize}
    \item capture catalogue ;
    \item capture filtres ;
    \item capture pagination.
\end{itemize}

\subsection{Génération des recommandations}

Le moteur de recommandation a été validé avec plusieurs projets utilisateurs.

Les tests ont confirmé :
\begin{itemize}
    \item la détection des compétences ;
    \item l’identification des catégories ;
    \item le calcul des scores ;
    \item le classement des cours.
\end{itemize}

\subsection{Exemple de recommandation}

Projet testé :

\begin{quote}
\textit{
"I want to create a complete website using React and Java"
}
\end{quote}

Résultats obtenus :
\begin{itemize}
    \item détection des compétences React et Java ;
    \item catégories Backend Development et Web Development ;
    \item génération de recommandations classées ;
    \item score explicable pour chaque résultat.
\end{itemize}

\vspace{0.5cm}

\begin{figure}[H]
    \centering

    \maybeincludegraphics[width=0.95\textwidth]{figures/recommendation-result.png}

    \fbox{
        \parbox[c][7cm][c]{0.9\textwidth}{
            \centering
            Placeholder :\\
            Capture des résultats de recommandation
        }
    }

    \caption{Résultats du moteur de recommandation}
\end{figure}

\subsection{Tests automatisés}

Le backend contient plusieurs tests automatisés permettant de vérifier le bon fonctionnement des composants critiques.

\begin{table}[H]
\centering
\caption{Tests automatisés principaux}
\begin{tabular}{|p{6cm}|p{7cm}|}
\hline
\textbf{Test} & \textbf{Objectif} \\
\hline
LoginAttemptServiceTests & Vérification du blocage après plusieurs tentatives échouées \\
\hline
RecommendationServiceTests & Validation de la normalisation et extraction des mots-clés \\
\hline
SkillBridgeApplicationTests & Vérification du chargement du contexte Spring Boot \\
\hline
\end{tabular}
\end{table}
\section{Validation du pipeline Big Data}

\subsection{Validation HDFS}

Le cluster HDFS a été validé avec :
\begin{itemize}
    \item deux DataNodes actifs ;
    \item réplication correcte des blocs ;
    \item absence de blocs manquants ;
    \item stockage des fichiers batch et streaming.
\end{itemize}

\vspace{0.3cm}

\textbf{Placeholder :}
\begin{itemize}
    \item capture dfsadmin -report ;
    \item capture DataNodes ;
    \item capture HDFS browser.
\end{itemize}

\subsection{Validation Sqoop}

Les imports batch ont été vérifiés via :
\begin{itemize}
    \item la présence des fichiers \texttt{\_SUCCESS} ;
    \item les fichiers \texttt{part-m-00000} ;
    \item les dossiers générés dans HDFS.
\end{itemize}

\vspace{0.3cm}

\textbf{Placeholder :}
\begin{itemize}
    \item capture import Sqoop ;
    \item capture dossiers HDFS.
\end{itemize}

\subsection{Validation Flume}

Les tests Flume ont confirmé :
\begin{itemize}
    \item la collecte des événements ;
    \item le transfert vers HDFS ;
    \item la création des fichiers events.*;
    \item le fonctionnement du pipeline streaming.
\end{itemize}

\subsection{Validation Hive}

Les requêtes Hive ont permis :
\begin{itemize}
    \item de lire les données HDFS ;
    \item d’analyser les événements ;
    \item d’exécuter des agrégations SQL ;
    \item de compter les événements utilisateurs.
\end{itemize}

\vspace{0.3cm}

\textbf{Placeholder :}
\begin{itemize}
    \item capture requêtes Hive ;
    \item capture résultats SQL.
\end{itemize}

\subsection{Validation MapReduce}

Le job \texttt{TopSearchKeywordsJob} a été exécuté avec succès.

Résultats observés :
\begin{itemize}
    \item extraction des mots-clés ;
    \item traitement distribué ;
    \item génération des statistiques ;
    \item écriture des résultats dans HDFS.
\end{itemize}

Les principaux mots-clés détectés sont :
\begin{itemize}
    \item react ;
    \item python ;
    \item machine ;
    \item spring ;
    \item frontend ;
    \item hadoop.
\end{itemize}

\vspace{0.3cm}

\textbf{Placeholder :}
\begin{itemize}
    \item capture MapReduce ;
    \item capture output keywords ;
    \item capture logs job.
\end{itemize}

\subsection{Validation HBase}

La table \texttt{course\_stats} a été générée et chargée avec succès.

Les statistiques stockées incluent :
\begin{itemize}
    \item nombre de clics ;
    \item sauvegardes ;
    \item progression moyenne ;
    \item informations de cours.
\end{itemize}

\vspace{0.3cm}

\textbf{Placeholder :}
\begin{itemize}
    \item capture scan HBase ;
    \item capture course\_stats ;
    \item capture shell HBase.
\end{itemize}

\section{Résultats obtenus}

\subsection{Résultats applicatifs}

Les principales fonctionnalités réalisées sont :
\begin{itemize}
    \item authentification sécurisée ;
    \item catalogue dynamique ;
    \item gestion des projets ;
    \item moteur de recommandation ;
    \item sauvegarde des cours ;
    \item suivi pédagogique ;
    \item dashboard administrateur.
\end{itemize}

\subsection{Résultats analytiques}

Le pipeline Big Data a permis :
\begin{itemize}
    \item l’ingestion des datasets ;
    \item le stockage distribué ;
    \item l’analyse des événements ;
    \item le traitement MapReduce ;
    \item la production d’indicateurs analytiques.
\end{itemize}

\subsection{Résumé des résultats}

\begin{table}[H]
    \centering
    \caption{Résumé des résultats du projet}

    \begin{tabular}{|p{6cm}|p{5cm}|}
        \hline
        \textbf{Composant} & \textbf{État} \\
        \hline

        Frontend React &
        Fonctionnel \\
        
        \hline

        Backend Spring Boot &
        Fonctionnel \\
        
        \hline

        Authentification JWT &
        Fonctionnelle \\
        
        \hline

        PostgreSQL / Supabase &
        Fonctionnel \\
        
        \hline

        Pipeline HDFS &
        Fonctionnel \\
        
        \hline

        Import Sqoop &
        Validé \\
        
        \hline

        Streaming Flume &
        Fonctionnel \\
        
        \hline

        Requêtes Hive &
        Fonctionnelles \\
        
        \hline

        Job MapReduce &
        Exécuté avec succès \\
        
        \hline

        HBase &
        Fonctionnel \\
        
        \hline
    \end{tabular}
\end{table}

\section{Discussion}

Le projet SkillBridge démontre la possibilité d’intégrer :
\begin{itemize}
    \item une application full-stack moderne ;
    \item un moteur de recommandation ;
    \item un pipeline Big Data complet ;
    \item des traitements analytiques distribués.
\end{itemize}

L’architecture mise en place présente plusieurs avantages :
\begin{itemize}
    \item séparation claire des responsabilités ;
    \item modularité ;
    \item évolutivité ;
    \item sécurité centralisée ;
    \item intégration analytique avancée.
\end{itemize}

Cependant, certaines limites subsistent :
\begin{itemize}
    \item pipeline partiellement manuel ;
    \item absence de traitement temps réel avancé ;
    \item moteur de recommandation rule-based ;
    \item dépendance aux scripts de démonstration.
\end{itemize}

\section{Conclusion}

Ce chapitre a présenté les différentes phases de validation réalisées sur le projet SkillBridge.

Les résultats obtenus démontrent :
\begin{itemize}
    \item le bon fonctionnement de l’application ;
    \item la stabilité du pipeline Big Data ;
    \item l’intégration réussie des technologies Hadoop ;
    \item la capacité du système à produire des recommandations et des analyses exploitables.
\end{itemize}

L’ensemble du projet constitue ainsi une plateforme full-stack et Big Data cohérente, modulaire et défendable académiquement.
*

\chapter{Validation académique}

\section{JEE et Spring Boot}

\begin{table}[H]
\centering
\caption{Concepts JEE appliqués}
\begin{tabular}{|p{5cm}|p{8cm}|}
\hline
\textbf{Concept} & \textbf{Implémentation} \\
\hline
Architecture en couches & Controllers / Services / Repositories / DTOs \\
\hline
API REST & Endpoints REST sécurisés \\
\hline
Persistance JPA & Entités relationnelles avec Spring Data JPA \\
\hline
Validation & Bean Validation et DTO validation \\
\hline
Transactions & Services transactionnels \\
\hline
\end{tabular}
\end{table}

\section{Spring Security}

\begin{table}[H]
\centering
\caption{Concepts Spring Security appliqués}
\begin{tabular}{|p{5cm}|p{8cm}|}
\hline
\textbf{Concept} & \textbf{Implémentation} \\
\hline
JWT & Authentification stateless \\
\hline
OAuth & Login Google et GitHub \\
\hline
Autorisation & Rôles USER / ADMIN \\
\hline
BCrypt & Hashage des mots de passe \\
\hline
Filtres de sécurité & JwtAuthenticationFilter \\
\hline
\end{tabular}
\end{table}

\section{Big Data}

\begin{table}[H]
\centering
\caption{Concepts Big Data appliqués}
\begin{tabular}{|p{5cm}|p{8cm}|}
\hline
\textbf{Concept} & \textbf{Implémentation} \\
\hline
Ingestion batch & Sqoop \\
\hline
Streaming & Flume \\
\hline
Stockage distribué & HDFS \\
\hline
SQL distribué & Hive \\
\hline
Traitement distribué & MapReduce \\
\hline
NoSQL & HBase \\
\hline
\end{tabular}
\end{table}

% ---------------------------------------------------------------------------
\chapter*{Conclusion générale}
\addcontentsline{toc}{chapter}{Conclusion générale}

Le projet \textbf{SkillBridge} avait pour objectif de concevoir une plateforme éducative intelligente capable de transformer une idée de projet en parcours d’apprentissage personnalisé.

Afin d’atteindre cet objectif, une architecture complète combinant :
\begin{itemize}
    \item développement full-stack ;
    \item sécurité applicative ;
    \item système de recommandation ;
    \item et pipeline Big Data ;
\end{itemize}
a été mise en place.

\vspace{0.5cm}

La plateforme développée permet :
\begin{itemize}
    \item la gestion des utilisateurs ;
    \item l’authentification sécurisée ;
    \item la consultation du catalogue de cours ;
    \item la création de projets ;
    \item la génération de recommandations personnalisées ;
    \item le suivi pédagogique ;
    \item ainsi que la supervision analytique via un dashboard administrateur.
\end{itemize}

Le moteur de recommandation constitue l’un des éléments centraux du projet. Celui-ci repose sur une approche déterministe et explicable basée sur :
\begin{itemize}
    \item l’extraction de mots-clés ;
    \item la détection des compétences ;
    \item l’identification des catégories ;
    \item et un système de scoring transparent.
\end{itemize}

Cette approche permet de produire des recommandations compréhensibles, traçables et facilement interprétables.

\vspace{0.5cm}

Le projet intègre également un pipeline Big Data complet utilisant plusieurs composants de l’écosystème Hadoop :
\begin{itemize}
    \item HDFS ;
    \item Sqoop ;
    \item Flume ;
    \item Hive ;
    \item MapReduce ;
    \item HBase.
\end{itemize}

Cette architecture permet :
\begin{itemize}
    \item l’ingestion de données massives ;
    \item le stockage distribué ;
    \item le traitement batch ;
    \item l’analyse des événements utilisateurs ;
    \item et la production d’indicateurs analytiques.
\end{itemize}

L’un des points importants du projet réside dans la séparation claire entre :
\begin{itemize}
    \item la couche transactionnelle temps réel ;
    \item et la couche analytique Big Data.
\end{itemize}

Cette organisation améliore :
\begin{itemize}
    \item les performances ;
    \item la modularité ;
    \item la maintenabilité ;
    \item et l’évolutivité du système.
\end{itemize}

\vspace{0.5cm}

Les validations réalisées ont permis de confirmer :
\begin{itemize}
    \item le bon fonctionnement de l’application ;
    \item la stabilité des services ;
    \item la communication entre les composants ;
    \item le bon fonctionnement du pipeline analytique ;
    \item ainsi que la cohérence des résultats générés.
\end{itemize}

Malgré les résultats obtenus, certaines limites subsistent :
\begin{itemize}
    \item le pipeline reste partiellement manuel ;
    \item le moteur de recommandation est principalement rule-based ;
    \item certains traitements analytiques pourraient être davantage automatisés ;
    \item et le système n’intègre pas encore de modèles avancés d’intelligence artificielle.
\end{itemize}

\vspace{0.5cm}

Plusieurs perspectives d’évolution peuvent être envisagées :
\begin{itemize}
    \item intégration de modèles NLP avancés ;
    \item utilisation d’embeddings sémantiques ;
    \item ajout de modèles de Machine Learning ;
    \item intégration de Spark ou Kafka ;
    \item automatisation complète du pipeline ;
    \item déploiement cloud distribué ;
    \item amélioration des dashboards analytiques ;
    \item personnalisation avancée des recommandations.
\end{itemize}

\vspace{0.5cm}

En conclusion, le projet SkillBridge constitue une plateforme cohérente combinant application web moderne et pipeline Big Data distribué. Il démontre l’intégration réussie de plusieurs technologies complémentaires dans une architecture modulaire et défendable académiquement.

Ce projet a également permis de mobiliser et consolider des compétences dans plusieurs domaines :
\begin{itemize}
    \item développement full-stack ;
    \item sécurité applicative ;
    \item gestion des bases de données ;
    \item traitement Big Data ;
    \item architecture logicielle ;
    \item et analyse de données distribuées.
\end{itemize}

\end{document}

